% !TeX root = ../../Book3.tex
% 纲要标题：### 12.6* 模的忠实平坦下降
% 纲要位置：工作记录/计划纲要.md，第 2327 行。
% 纲要细目（写作范围）：
% * 双重基变换及余循环条件
% * 忠实平坦下降的模论形式
% * 严格区分带下降数据与只有基变换后对象
% ---

\OptionalSection{模的忠实平坦下降}{b3:sec:1206}

给定忠实平坦 \(R\) 代数 \(S\) 上的模 \(N\)，什么时候能写成 \(S\otimes_RM\)？仅有 \(N\) 还不够；它在两个 \(S\) 因子上的拉回之间必须有指定的相容识别。本节把这些数据写清楚，再用一个等化子实际构造 \(M\)。


\subsection{双重基变换与余循环条件}
\label{b3:subsec:1206:01}

固定忠实平坦环同态 \(R\to S\)，令 \(T=S\otimes_RS\)。同一个 \(S\) 模 \(N\) 有两种变基：
\[
N_1=N\otimes_RS,\qquad N_2=S\otimes_RN.
\]
在前者上，\((a\otimes b)(n\otimes s)=an\otimes bs\)；在后者上，\((a\otimes b)(s\otimes n)=as\otimes bn\)。这两种 \(T\) 模结构一般不能忽略下标而混写。

\begin{definition}\label{b3:def:module-descent-datum}
设 \(R\to S\) 为忠实平坦交换环同态，\(N\) 为 \(S\) 模，\(T=S\otimes_RS\)。赋予 \(N_1=N\otimes_RS\) 及 \(N_2=S\otimes_RN\) 自然的 \(T\) 模结构，即两个 \(S\) 因子分别作用在相应位置。若 \(T\) 线性同构
\(\theta:N_1\to N_2\) 满足下面的\term[yuxunhuan-tiaojian]{余循环条件}[cocycle condition]，就称它为 \(N\) 上的\term[xiajiang-shuju]{下降数据}[descent datum]。具体地，把数据沿三个嵌入
\[
T\longrightarrow S\otimes_RS\otimes_RS,
\quad a\otimes b\longmapsto
a\otimes b\otimes1,\ a\otimes1\otimes b,\ 1\otimes a\otimes b
\]
变基，分别得到从第一个位置移到第二个位置的 \(\theta_{12}\)、从第一个移到第三个的 \(\theta_{13}\)、从第二个移到第三个的 \(\theta_{23}\)。要求
\[
\theta_{23}\theta_{12}=\theta_{13}.
\]
把三重张量积中第 \(i\) 个 \(S\) 因子换成 \(N\) 所得的模记为 \(N^{(i)}\)，其中 \(i=1,2,3\)。余循环条件就是下面三个变基同构构成的图交换：
\[
\begin{tikzcd}[column sep=large,row sep=large]
N^{(1)}\arrow[r,"\theta_{12}"]\arrow[dr,"\theta_{13}"']
&N^{(2)}\arrow[d,"\theta_{23}"]\\
&N^{(3)}
\end{tikzcd}
\]
带下降数据的模之间的态射，是与这些同构相容的 \(S\) 线性映射。
\end{definition}

若 \(N=S\otimes_RM\)，标准下降数据为
\[
\theta\bigl((s\otimes m)\otimes t\bigr)
 =s\otimes(t\otimes m).
\]
它的逆把右边送回左边。在三重张量积上，两种复合都把位于第一个 \(S\) 因子旁的 \(m\) 移到第三个因子旁，系数次序不变，所以满足余循环条件。

\ProseParagraphBreak
为把一般条件用于计算，令
\[
\rho:N\longrightarrow S\otimes_RN,\qquad
\rho(n)=\theta(n\otimes1),\qquad
\mu(s\otimes n)=sn.
\]
\(\rho\) 是 \(R\) 线性映射，满足
\begin{equation}\label{b3:eq:descent-coaction}
\begin{split}
\rho(sn)&=(s\otimes1)\rho(n),\qquad \mu\rho(n)=n,\\
(1\otimes\rho)\rho(n)&=\sum_i s_i\otimes1\otimes n_i
\quad\text{若 }\rho(n)=\sum_i s_i\otimes n_i.
\end{split}
\end{equation}
第一式是 \(T\) 线性。第三式就是把余循环条件用于 \(n\otimes1\otimes1\)：先移动到第二个位置，再移动到第三个，得到左边；直接移到第三个得到右边。第二式来自对角变基：沿 \(S\otimes_RS\to S\)、\(s\otimes t\mapsto st\) 拉回 \(\theta\)，得到 \(N\) 的自同构 \(u=\mu\rho\)。再沿全对角映射
\[
S\otimes_RS\otimes_RS\longrightarrow S,
\qquad s_1\otimes s_2\otimes s_3\longmapsto s_1s_2s_3
\]
变基，三个 \(N^{(i)}\) 都自然认同于 \(N\)，而 \(\theta_{12},\theta_{13},\theta_{23}\) 都变为同一个 \(u\)，因为三个 \(T\) 嵌入与全对角映射的复合都是乘法映射 \(T\to S\)。于是余循环条件成为 \(u^2=u\)。\(u\) 可逆，约去一个 \(u\) 得 \(u=\operatorname{id}_N\)，即 \(\mu\rho=\operatorname{id}_N\)。原同构可由 \(\rho\) 恢复为 \(\theta(n\otimes t)=(1\otimes t)\rho(n)\)。

\subsection{忠实平坦下降的模论形式}
\label{b3:subsec:1206:02}

先解释忠实平坦性为何能够恢复本来就在 \(R\) 上的模。

\begin{lemma}\label{b3:lem:faithful-flat-equalizer}
设 \(R\to S\) 为忠实平坦交换环同态，\(M\) 为任意 \(R\) 模。则下列序列正合：
\[
0\longrightarrow M\xrightarrow{m\mapsto1\otimes m}S\otimes_RM
\xrightarrow{\mathrm{d}}S\otimes_RS\otimes_RM,
\]
其中 \(\mathrm{d}(s\otimes m)=1\otimes s\otimes m-s\otimes1\otimes m\)。
\end{lemma}
\begin{proof}
由\cref{b3:prop:faithfully-flat-detection}，可以先与 \(S\) 张量再检查。新序列的首个映射将 \(a\otimes m\) 送到 \(a\otimes1\otimes m\)，乘合前两个 \(S\) 因子给出左逆，故单射。

若 \(z=\sum a_i\otimes b_i\otimes m_i\) 在后一个映射的核中，则
\[
\sum_i a_i\otimes1\otimes b_i\otimes m_i
=\sum_i a_i\otimes b_i\otimes1\otimes m_i.
\]
乘合前两个 \(S\) 因子，得到
\(z=\sum_i a_ib_i\otimes1\otimes m_i\)，恰好属于前一个映射的像。相反，该像代入差映射显然为零。变基后的正合性因而成立，并反映回原序列。
\end{proof}

\begin{theorem}[忠实平坦模下降]\label{b3:thm:faithfully-flat-module-descent}
设 \(R\to S\) 为忠实平坦交换环同态，\(N\) 为带下降数据 \(\theta:N\otimes_RS\to S\otimes_RN\) 的 \(S\) 模，记 \(\rho(n)=\theta(n\otimes1)\)，并令
\[
M=\bigl\{n\in N:\rho(n)=1\otimes n\bigr\}.
\]
则 \(M\) 是 \(R\) 子模，且乘法映射
\[
\varphi:S\otimes_RM\longrightarrow N,\qquad s\otimes m\longmapsto sm
\]
为保下降数据的同构。这个构造与基变换建立 \(R\) 模范畴和带下降数据的 \(S\) 模范畴之间的等价。若 \(N\) 有限展示，则所得 \(M\) 有限展示。
\end{theorem}
\begin{proof}
记 \(\eta:N\to S\otimes_RN\)、\(\eta(n)=1\otimes n\)。\(\rho\) 与 \(\eta\) 都是 \(R\) 线性映射，且 \(M=\ker(\rho-\eta)\)，所以有左端正合列
\[
0\longrightarrow M\longrightarrow N
\xrightarrow{\ \rho-\eta\ }S\otimes_RN.
\]
与 \(S\) 张量后，由平坦性仍有左端正合列
\[
0\longrightarrow S\otimes_RM\longrightarrow S\otimes_RN
\xrightarrow{\ 1\otimes(\rho-\eta)\ }S\otimes_RS\otimes_RN.
\]
因此可将 \(S\otimes_RM\) 认同于 \(S\otimes_RN\) 中下列两个映射的等化子：
\[
s\otimes n\longmapsto s\otimes\rho(n),\qquad
s\otimes n\longmapsto s\otimes1\otimes n.
\]
\eqref{b3:eq:descent-coaction}的第三式说明 \(\rho(n)\) 落在这个等化子中，故 \(\rho\) 给出 \(N\to S\otimes_RM\) 的映射。第二式说明 \(\varphi\rho=1_N\)。

反过来，若 \(z=\sum_i s_i\otimes n_i\) 属于上述等化子，则
\(\sum_i s_i\otimes\rho(n_i)=\sum_i s_i\otimes1\otimes n_i\)。乘合前两个因子，并用第一式，得到
\[
\rho\left(\sum_i s_in_i\right)=\sum_i s_i\otimes n_i=z.
\]
所以 \(\rho\varphi\) 也是恒等映射，\(\varphi\) 为同构。对 \(m\in M\)，\(\rho(sm)=s\otimes m\)，这正是标准数据在 \(s\otimes m\) 上的公式，故同构与原 \(\theta\) 相容。

若 \(u:N\to N'\) 与数据相容，则 \(u\) 将 \(\rho(n)=1\otimes n\) 的元素送入对应的 \(M'\)，因而限制为 \(R\) 线性映射 \(M\to M'\)。它在 \(S\) 上的扩张通过 \(\varphi\) 恢复原 \(u\)。若从 \(R\) 模 \(M_0\) 出发，\cref{b3:lem:faithful-flat-equalizer}恰说明标准数据的等化子是 \(M_0\) 的自然像，故往返构造为自然同构。这证明范畴等价的对象与态射两部分。最后，\(N\cong S\otimes_RM\)，有限展示由\cref{b3:prop:faithfully-flat-detection}第四项下降。
\end{proof}

等化子用相容性选出能够来自原环的元素，平坦性使这个核的构造与基变换相容；忠实性则保证从原环出发的对象与态射能完整恢复。

\subsection{下降数据与基变换后对象}
\label{b3:subsec:1206:03}

仅给出变基后的对象，还没有指定其中哪些元素应从原环而来。例如 \(R=\mathbb R\)、\(S=\mathbb C\)，在 \(N=\mathbb C\) 上可考虑两种半线性对合
\[
\sigma_+(z)=\bar z,\qquad \sigma_-(z)=-\bar z.
\]
它们都满足 \(\sigma(cz)=\bar c\cdot\sigma(z)\)、\(\sigma^2=1\)，分别选出不动实子空间 \(\mathbb R\) 与 \(i\mathbb R\)。这里
\(\mathbb C\otimes_{\mathbb R}\mathbb C\cong\mathbb C\times\mathbb C\)，映射为 \(z\otimes w\mapsto(zw,z\bar w)\)。这两个分量分别对应恒等嵌入与共轭嵌入。在第一个分量上，\(N_1,N_2\) 都识别为 \(N\)，\(\theta\) 是恒等映射；在第二个分量上，\(N_1\) 为 \(N\)，\(N_2\) 为共轭标量作用的模 \(\overline N\)，其中 \(c\cdot n=\bar c n\)。此处 \(\theta\) 由 \(n\mapsto\sigma(n)\) 给出：半线性等式 \(\sigma(cn)=\bar c\sigma(n)\) 正是从 \(N\) 到 \(\overline N\) 的复线性。因而，两种对合都给出前面定义中的 \(T\) 线性同构。

\ProseParagraphBreak
三重张量积的分量可以用两次相对嵌入 \(g,h\in\{1,\text{共轭}\}\) 标记。\(\theta_{23}\theta_{12}=\theta_{13}\) 在这些分量上要求连续两次相对识别等于复合嵌入的识别；涉及恒等嵌入的条件自动成立，唯一剩下的条件就是 \(\sigma^2=1\)。等化子条件在恒等分量上自动成立，在共轭分量上则为 \(\sigma(n)=n\)，故恢复的实模恰为 \(N^\sigma\)。在每种情况下，乘法映射 \(\mathbb C\otimes_{\mathbb R}N^\sigma\to N\) 都是同构。这两组数据也同构：复线性映射 \(q(z)=iz\) 满足
\[
\sigma_-\bigl(q(z)\bigr)=i\bar z=q\bigl(\sigma_+(z)\bigr),
\]
因而与下降数据相容，并将 \(\mathbb R\) 同构地送到 \(i\mathbb R\)。这个例子说明，\(N\) 内具体的固定子模及其乘法识别，需要由下降数据指定；所得实向量空间的同构类型在这里相同。

\ProseParagraphBreak
开覆盖的例子把下降数据写成换基矩阵。设 \(R\) 为交换环，\(r\) 为非负整数，且 \(D(f_1),\ldots,D(f_q)\) 组成 \(\operatorname{Spec}R\) 的有限主开覆盖。在各开集上取自由模 \(N_i=R_{f_i}^r\)，令
\[
S=\prod_{i=1}^qR_{f_i},\qquad
N=\prod_{i=1}^qN_i\cong S^r.
\]
上一节已证明 \(R\to S\) 忠实平坦，且 \(S\otimes_RS\cong\prod_{i,j}R_{f_if_j}\)。在重叠处指定矩阵
\[
g_{ij}\in\operatorname{GL}_r(R_{f_if_j}),
\qquad g_{jk}g_{ij}=g_{ik}\quad\text{于 }R_{f_if_jf_k}.
\]
这些 \(r\) 阶可逆矩阵给出相容同构 \(\theta_{ij}\)，从而给出 \(N\) 的下降数据。所得等化子 \(M\) 由重叠处相符的局部元素组组成，下降同构在第 \(i\) 个分量上给出 \(M_{f_i}\cong N_i=R_{f_i}^r\)。因此 \(M\) 在主开覆盖上有限局部自由，由\cref{b3:thm:finite-presentation-flat-projective}为有限生成投射模。这些矩阵未必能由各个开集上的换基同时消去，所以不一定得到全局自由模。

\ProseParagraphBreak
余循环条件保证局部识别相容，等化子将它们还原成原环上的模；有限展示性也随之下降，使有限组生成元与关系的描述得以保留。
